\newpage
\section{Conclusion et Perspectives}
\label{sec:conclusion}

La tâche 3.2 posait deux questions de recherche sur l'explicabilité des modèles prédictifs profonds de séries temporelles, appliquées à la prévision du niveau des nappes. L'instruction de ces questions a fait apparaître un préalable qui n'était pas anticipé et qui a orienté une part importante du travail : la donnée hydrologique française, quoique publique et de bonne qualité, n'était pas exploitable en l'état pour l'apprentissage automatique, faute de jointure entre les compartiments, de consolidation de l'historique et d'indicateurs normalisés.

Le premier acquis est donc la levée de ce verrou. L'entrepôt répondant à QT1 agrège quatre sources ouvertes en une chaîne idempotente, reprenable et déployable, expose des chroniques où le niveau et son forçage climatique figurent sur la même ligne au même jour, et produit une famille d'indicateurs standardisés dont la justesse a été vérifiée contre un critère fixé avant la mesure. La plateforme répondant à QT2 y ajoute l'exploration cartographique, un laboratoire de modèles mécanistiques calibrés selon les critères STOWA et paramétrés par famille d'aquifère, un catalogue d'architectures de prévision profonde sous traçabilité d'expériences, et les modules d'explication. Ce que cette chaîne change se mesure à une chose : obtenir soixante ans de niveaux de nappe joints à leur climat demandait plusieurs semaines de préparation, c'est aujourd'hui une requête.

Le deuxième acquis porte sur QR1. Cinq familles d'attribution sont mises en œuvre, dont quatre accessibles dans l'interface, non pour l'exhaustivité mais pour la confrontation : leur convergence rend une explication crédible, leur divergence signale une explication fragile, information qu'une méthode unique aurait masquée. La restitution projette l'attribution sur le plan variable-temps, forme sous laquelle un hydrogéologue peut la confronter à sa connaissance des temps de réponse de l'aquifère. La question de savoir si ce dispositif permet effectivement aux experts de contrôler la qualité décisionnelle d'un modèle reste ouverte : elle appelle un protocole avec utilisateurs qui n'a pas été conduit.

Le troisième acquis porte sur QR2, et il est double. Il a d'abord fallu constater que la formulation standard du problème contrefactuel ne convient pas au domaine : une perturbation libre des entrées produit des scénarios valides au sens de l'optimisation et physiquement absurdes, et traiter la plausibilité comme un terme de régularisation ne suffit pas, parce qu'un scénario qui viole un peu les relations physiques est rejeté aussi sûrement qu'un scénario qui les viole beaucoup. Deux méthodes de natures différentes ont donc été implémentées, l'une par paramétrisation à faible dimension d'un régime climatique, l'autre par recherche d'un analogue réellement observé, et un test de cohérence croisée contre un modèle mécanistique calibré indépendamment permet de rejeter un scénario sortant du domaine de validité. Ce test est le seul contrôle disponible sur un point où, par construction, aucune observation n'existe, et il est indépendant de la méthode de génération, donc applicable à toute méthode qu'on intégrerait au dispositif. Ces développements sont livrés en bibliothèque et exposés par l'API, mais la page qui devait les mettre à portée d'un hydrogéologue n'est branchée dans aucune route de l'application : ils ne s'exécutent aujourd'hui que par l'API ou par un carnet de calcul.

Ce qui n'est pas acquis a été énoncé en section~\ref{sec:benchmark} et mérite d'être répété ici : ni le classement des architectures de prévision sur le cas hydrogéologique, ni la comparaison quantitative des familles de contrefactuels selon les dimensions retenues n'ont été mesurés. Le livrable de cette partie de la tâche est le protocole qui rend ces mesures possibles et comparables, non leur résultat.

\subsection*{Perspectives}

Cinq chantiers se dégagent, par ordre de priorité.

Le premier est correctif et conditionne les autres. Le forçage d'évapotranspiration de la chaîne des stations doit être remplacé par l'évapotranspiration de référence, ce qui suppose de transporter les extrêmes journaliers de température jusqu'à cette chaîne, de recalculer la grandeur au pas journalier et de recalibrer les modèles existants. Tant que ce n'est pas fait, les paramètres des modèles mécanistiques n'ont pas d'interprétation physique directe, et aucune explication ne peut s'appuyer sur eux pour justifier une prévision.

Le deuxième est le branchement de l'interface contrefactuelle, sans lequel aucune évaluation associant des hydrogéologues n'est possible. La page et ses composants existent, l'API répond, seule la déclaration de route manque.

Le troisième est l'exécution du protocole de la section~\ref{sec:methodo}, sur un échantillon stratifié par famille d'aquifère, avec la référence naïve incluse et le traitement statistique par diagramme de différence critique. C'est de cette campagne que dépend la réponse à QR2, et l'outillage nécessaire est en place.

Le quatrième est l'extension du cadre à la prévision probabiliste. Les méthodes mises en œuvre expliquent une prévision ponctuelle, alors que la décision opérationnelle porte sur un risque de franchissement de seuil, donc sur une distribution. Des travaux récents formulent la saillance contrefactuelle pour des modèles de régression probabiliste \cite{raman2023information}, et le catalogue d'architectures de la plateforme comporte déjà des modèles produisant des quantiles, ce qui rend cette extension accessible sans reconstruction du socle.

Le cinquième est l'adaptation des configurations d'apprentissage par famille d'aquifère. La profondeur d'historique en entrée, l'horizon de prévision pertinent et le pas d'agrégation devraient être choisis selon la nature de la nappe, comme ils le sont déjà côté mécanistique. C'est vraisemblablement le premier gain de performance à aller chercher, et il conditionne la qualité des modèles sur lesquels porteront les explications.

Trois sujets identifiés avec les hydrogéologues du BRGM restent par ailleurs ouverts et dépassaient le périmètre de la tâche : l'étude des échanges nappe-rivière et du soutien d'étiage, la caractérisation de la double porosité, et la définition d'un indicateur de situation pour les nappes captives, l'indicateur standardisé existant n'ayant été conçu que pour les nappes libres. Ils sont consignés ici parce qu'ils viennent des experts du domaine et constituent des points d'entrée pour la suite.

L'ensemble de ces développements est public sous licence MIT~\cite{depot_entrepot, depot_plateforme}, documenté et déployable depuis une machine vierge, de sorte que les partenariats à venir entre le BRGM et le LIFAT, notamment dans le cadre du projet ANR AIDA, disposent d'un socle qu'ils n'auront pas à reconstruire.
